% F1: das transformações dos dados aos braços. Diagrama em TikZ (nenhum dado de resultado).
% Fica na raiz de paper/ porque o Overleaf do artigo não tem pastas.
\vspace{2mm}
\noindent\begin{minipage}{\columnwidth}
\centering\small
\begin{tikzpicture}[
  box/.style={draw, rounded corners=2pt, align=center, inner sep=4pt, text width=0.88\columnwidth},
  arr/.style={-{Stealth[length=2mm]}}]
\node[box] (raw) at (0,0) {Eventos brutos (\nEvents{} eventos)\\ lidos pelo gabarito};
\node[box] (act) at (0,-2.5) {Tabela de ações (\nActions{} ações)\\ lida por \armEventsInPrompt{} e \armVector{}};
\node[box] (grafo) at (0,-5.0) {Grafo Neo4j e súmula agregada\\ lidos por \armStatsInPrompt{}, \armTextToCypher{} e \armGraphTools{}};
\draw[arr] (raw) -- node[right, xshift=1mm] {SPADL, zonas, posse e xT} (act);
\draw[arr] (act) -- node[right, xshift=1mm] {carga no banco} (grafo);
\end{tikzpicture}
\captionof{figure}{Transformações dos dados e braços que leem cada camada.}\label{fig:dados}
\end{minipage}\par\smallskip
